\begin{frame}{例子：一次多轮对话的真实链路并不止于模型前向}
  \Core{请求生命周期要覆盖从输入组织到状态回收的完整过程，而不是只盯住模型执行。}
  \begin{block}{进入前的准备}
    chat template、tokenizer 和 prompt 拼接本身就会引入时间与长度变化。
  \end{block}
  \pause
  \begin{block}{进入后的主路径}
    入队、调度、prefill、decode、输出组装会在不同对象之间多次转移控制权。
  \end{block}
  \pause
  \begin{block}{完成后的尾部动作}
    请求结束不代表系统工作结束，状态释放、统计更新和资源回收同样属于生命周期的一部分。
  \end{block}
\end{frame}

\begin{frame}{例子：chat template 为什么会影响生命周期分析}
  \Core{如果输入在真正进入 engine 前已经发生了长度膨胀，那么后面的阶段观察也会被一并放大。}
  \begin{block}{用户看到的是一条自然语言消息}
    但系统里常常先被改写为带角色标记、特殊 token 和 generation prompt 的模板串。
  \end{block}
  \pause
  \begin{block}{直接后果}
    prompt\_len 变化会影响 prefill 代价、KV 初始规模和后续排队行为。
  \end{block}
  \pause
  \begin{block}{教学意义}
    生命周期不是从“用户文本”开始算，而是从“真正进入系统的 token 序列”开始算。
  \end{block}
\end{frame}

\begin{frame}{代码例子：多轮对话在入口处如何被组织}
  \Core{生命周期课程里的代码例子要回答：输入是怎样变成请求对象的。}
  \begin{block}{入口片段}
    {\ttfamily\small
    \detokenize{prompts = [tokenizer.apply_chat_template(...)]}\par
    \detokenize{outputs = llm.generate(prompts, sampling_params)}\par
    }
  \end{block}
  \pause
  \begin{block}{课堂应追问}
    `generate()` 接到的是原始文本、token 序列，还是已经带运行时元数据的请求对象？
  \end{block}
  \pause
  \begin{block}{系统含义}
    入口处的数据形态，决定了后面 state object 该如何承载长度、状态和 block table。
  \end{block}
\end{frame}

\begin{frame}{代码例子：Sequence 为什么是生命周期里的关键对象}
  \Core{如果不能说清 \texttt{Sequence} 承载了哪些字段，就很难讲明请求是如何跨阶段存活的。}
  \begin{block}{课堂应关注的字段}
    {\ttfamily\small
    \detokenize{status}\par
    \detokenize{num_cached_tokens}\par
    \detokenize{block_table}\par
    \detokenize{last_block_num_tokens}\par
    }
  \end{block}
  \pause
  \begin{block}{对象视角}
    这些字段把“请求内容、调度状态、KV 元数据”压在了同一个生命周期对象上。
  \end{block}
  \pause
  \begin{block}{结论}
    一次请求之所以是系统问题，关键就在于它不是一次函数调用，而是一个持续演化的状态对象。
  \end{block}
\end{frame}

\begin{frame}{例子：prefill 和 decode 为什么应该拆开看}
  \Core{生命周期讲义里最重要的区分，就是 prefill 与 decode 的资源特征不同。}
  \begin{block}{prefill}
    一次性处理更长输入，计算更密集，也更容易直接抬高 TTFT。
  \end{block}
  \pause
  \begin{block}{decode}
    每轮 token 较少，但会跨更多轮持续存在，更容易与调度和状态驻留绑定在一起。
  \end{block}
  \pause
  \begin{block}{误区}
    如果把二者都写成“模型生成阶段”，就会把关键的性能来源混在一起。
  \end{block}
\end{frame}

\begin{frame}{课堂练习：用时序图复述一次请求}
  \Core{课程练习的目的不是画复杂图，而是检验你是否知道控制权何时转移。}
  \begin{block}{练习一}
    在时序图中标出：请求进入、prefill 开始、decode 开始、输出完成、状态释放五个节点。
  \end{block}
  \pause
  \begin{block}{练习二}
    说明其中哪一个节点最适合插第一条日志，为什么。
  \end{block}
  \pause
  \begin{block}{练习三}
    用一句话回答：为什么生命周期图是后续调度和 KV 讨论的公共前提？
  \end{block}
\end{frame}
